\section{Conclusiones}

El modelo físico mostró que, por eje, la planta es un doble integrador que no es estable en lazo
abierto. Por eso no aplicaban ni las reglas SIMC ni un PI, y la identificación tuvo que hacerse en
lazo cerrado. El modelo lineal es válido hasta unos \SI{3}{\degree} con cambios lentos de
inclinación, pero solo hasta \SI{0.5}{\degree} con cambios rápidos, por el término centrífugo.

Con el modelo de diseño, el controlador PD con integración condicional cumplió los requisitos en
MIL ($t_s = \SI{1.30}{\second}$, $M_p = \SI{4.4}{\percent}$, $e_{ss} = \SI{0.18}{\milli\metre}$),
mientras que el PID clásico llegó a \SI{38.8}{\percent} de sobreimpulso. La saturación de
\SI{10}{\degree} que impone el mecanismo no alcanza para frenar un empujón de
\SI{420}{\milli\metre\per\second} antes del borde.

El código generado con Embedded Coder fue equivalente al modelo en SIL y en PIL sobre la ESP32,
con igualdad exacta en todas las muestras y \SI{31.3}{\micro\second} de cálculo como máximo. Para
el PIL hubo que construir una ruta propia, porque el paquete de soporte de Arduino no lo habilita
para la ESP32. Estos resultados verifican el código, no el desempeño sobre la planta.

La identificación sobre el prototipo dio una ganancia efectiva de \numrange{0.12}{0.28}, un
retardo de unos \SI{0.3}{\second} y una zona muerta de alrededor de \SI{1}{\degree}, atribuibles
al mecanismo. El modelo de diseño no representa bien el prototipo, y eso explica que en la planta
real no se cumplieran R1 ni R3: los escalones no se establecieron en \SI{4.8}{\second} y
terminaron a entre \SI{9}{} y \SI{80}{\milli\metre} de la referencia, aunque sin sobreimpulso. El
seguimiento de un círculo de \SI{45}{\milli\metre} sí funcionó, con un error mediano de
\SI{30}{\milli\metre}.

Las limitaciones principales del trabajo son tres. La implementación física usó un firmware
escrito a mano con la misma ley, y no el código generado. El rediseño con el modelo identificado
solo se evaluó en el prototipo. Y la última versión del firmware, con el integrador corregido y el
nivel reajustado, no se alcanzó a probar. Para mejorar el desempeño hace falta, antes que otro
controlador, reducir la holgura y el retardo del mecanismo, sobre todo en el eje Y.
